\paragraph{}Niniejsza sekcja prezentuje wyniki ewaluacji modeli rozpoznawania twarzy na testowym zbiorze CASIA-WebFace, ze szczególnym uwzględnieniem ich odporności na okluzję. Wartości Rank-1/Rank-3 pochodzą z zapisów \texttt{score.txt} wygenerowanych przez pipeline FAISS w scenariuszu okluzji 20px. Wyniki zawierają porównanie dokładności Rank-k, analizę wpływu syntetycznej okluzji oraz porównanie z modelami SOTA.

\subsection{Porównanie z modelami SOTA}

\paragraph{}W celu ustalenia benchmarku, poddaliśmy testowaniu trzy modele stanowiące stan techniki w rozpoznawaniu twarzy: ArcFace Large (buffalo\_l), ArcFace Small (buffalo\_s) oraz VGGFace (model legacy). Wszystkie modele testowano w identycznym scenariuszu okluzji (syntetyczna okluzja 20 pikseli na poziomie oczu). Wyniki zawarto w Tabeli~\ref{tab:sota_baseline}.

\begin{table}[H]
\centering
\small
\begin{tabular}{lcc}
\toprule
\textbf{Model} & \textbf{Rank-1 (\%)} & \textbf{Rank-3 (\%)} \\
\midrule
ArcFace Large (buffalo\_l) & 94.81 & 95.51 \\
ArcFace Small (buffalo\_s) & 86.65 & 90.68 \\
VGGFace (Legacy) & 39.64 & 57.01 \\
\bottomrule
\end{tabular}
\caption{Wyniki modeli SOTA na zbiorze testowym z syntetyczną okluzją 20px.}
\label{tab:sota_baseline}
\end{table}

\paragraph{}ArcFace Large wyznacza górny limit wydajności (94.81\%), przy czym model wykorzystuje masywny dataset treningowy (MS1MV2 z $\sim$5.8M obrazów). ArcFace Small (86.65\%) stanowi punkt odniesienia dla efektywnych modeli przemysłowych. VGGFace (39.64\%) pokazuje drastyczną degradację, podkreślając znaczenie nowoczesnych funkcji straty opartych na marginesach (ArcFace) w porównaniu do starszych podejść opartych na softmax.

\subsection{Dokładność modeli badawczych z okluzją}

\paragraph{}Modele własne obecne w repozytorium obejmują wariant bazowy, wariant z CBAM w blokach rezydualnych oraz wariant z gałęzią Transformer. Wyniki Rank-1 i Rank-3 przedstawiono w Tabeli~\ref{tab:results_custom}.

\begin{table}[H]
\centering
\small
\begin{tabular}{lccc}
\toprule
\textbf{Model} & \textbf{Rank-1 (\%)} & \textbf{Rank-3 (\%)} & \textbf{Delta vs Baseline} \\
\midrule
Baseline & 87.11 & 90.56 & -- \\
CBAM\_Block & 90.16 & 91.88 & +3.05 pp \\
Transformers & 85.00 & 89.02 & -2.11 pp \\
\bottomrule
\end{tabular}
\caption{Dokładność modeli badawczych na zbiorze testowym z syntetyczną okluzją 20px.}
\label{tab:results_custom}
\end{table}

\paragraph{}Najlepsze wyniki wśród modeli badawczych osiągnął wariant CBAM\_Block (90.16\%), przekraczając baseline o 3.05 pp. Wariant Transformer uzyskał niższy wynik niż baseline.

\subsection{Analiza degradacji wydajności}

\paragraph{}Porównanie wyników modeli w scenariuszu okluzji, nieznaczenie, ale potwierdza, że rozproszona uwaga (CBAM\_Block) zwiększa odporność na zasłonięcie kluczowych regionów. Model ArcFace Large osiąga 94.81\% w warunkach okluzji 20px, podczas gdy najlepszy wariant badawczy CBAM\_Block uzyskuje 90.16\%.

\subsection{Ranking modeli i luka Rank-3}
\paragraph{}W celu analizy zachowania top-k zestawiono różnice Rank-3 vs Rank-1 (im mniejsza luka, tym większa pewność top-1). Najmniejszą lukę uzyskano dla ArcFace Large i CBAM\_Block, co wskazuje na stabilne dopasowania wśród najbliższych sąsiadów.

\begin{table}[H]
\centering
\small
\begin{tabular}{lccc}
\toprule
\textbf{Model} & \textbf{Rank-1 (\%)} & \textbf{Rank-3 (\%)} & \textbf{Luka Rank-3 - Rank-1} \\
\midrule
ArcFace Large (buffalo\_l) & 94.81 & 95.51 & 0.70 \\
CBAM\_Block & 90.16 & 91.88 & 1.72 \\
Baseline & 87.11 & 90.56 & 3.45 \\
ArcFace Small (buffalo\_s) & 86.65 & 90.68 & 4.03 \\
Transformers & 85.00 & 89.02 & 4.02 \\
VGGFace (Legacy) & 39.64 & 57.01 & 17.37 \\
\bottomrule
\end{tabular}
\caption{Luka pomiędzy Rank-3 i Rank-1 dla wszystkich modeli (mniejsza luka = większa pewność top-1).}
\label{tab:rank_gap}
\end{table}